\section{Background and Related Work}
\label{sec:literature}

This section first summarises the gains that situation awareness can bring to
next-generation access networks, and then discusses the simulation tooling
that is required to study these gains in a site-specific manner. The latter
discussion directly motivates the choice of ns-3/5G~LENA combined with Sionna
RT adopted in this paper.

\subsection{Benefits of Situation Awareness}

Situation awareness refers to the capability of antenna systems to adapt their
configuration based on contextual information such as UE position, mobility,
environment geometry, obstacles, and temporal channel variation. Recent
studies show that incorporating such awareness leads to substantial gains in
performance, reliability, and efficiency, which can be grouped into four
families.

\subsubsection{Beamforming and Beam Management}
Surveys of AI-enabled multimodal beamforming indicate that combining non-RF
sensors (e.g., LiDAR, radar, GPS) with RF data improves directional accuracy
and reduces the overhead of beam sweeping and alignment, particularly in
mmWave bands where misalignment and rapid channel variation degrade
performance~\cite{ref:multimodal_bf}. Positioning-aided beamforming for mmWave
communications has specifically shown how location information can relax the
complexity of beam training and narrow the search space for optimal beams.

\subsubsection{Mobility and Handover Management}
Context-aware and predictive handover schemes that exploit UE speed,
direction, hysteresis, trajectory history, and signal-quality metrics reduce
handover failures, ping-pong effects, and latency. Memory-full context-aware
predictive mobility management in dual-connectivity 5G networks uses past UE
trajectory and signal measurements to anticipate handover events and reduce
signalling overhead~\cite{ref:predictive_handover}. Surveys on AI-aided
mobility and handover management confirm that machine-learning-based schemes
can sustain seamless connectivity even in ultra-dense deployments and
high-mobility regimes.

\subsubsection{Reliability and Link Stability}
Awareness of obstacles and environmental geometry enables dynamic beam
switching or alternate-path selection, mitigating the impact of blockages
especially in higher-frequency bands. Future networks are expected to rely
heavily on integrated sensing (and positioning) capabilities to improve link
robustness under dynamic conditions~\cite{ref:positioning_6g}.

\subsubsection{Resource Efficiency and Latency Reduction}
Predictive mobility and proactive beam adjustment avoid unnecessary sweeping,
measurement reporting, and control overhead. Deep-learning-based predictive
handover algorithms have been shown to reduce interruption time and improve
resource utilisation; in mobile edge-computing systems, joint optimisation of
connectivity and computation likewise benefits from context-aware decision
rules~\cite{ref:predictive_handover}.

While the literature confirms substantial improvement potential, it also
emphasises trade-offs: acquiring accurate context (position, environment
mapping) has its own overhead, prediction errors can degrade rather than
improve performance, and privacy and hardware constraints must be addressed.

\subsection{Simulation Tooling for ISAC}

Studying situation-aware ISAC requires a simulation environment that captures
both end-to-end networking behaviour and site-specific propagation. Two
simulator families are commonly used as starting points.

\textbf{Simu5G on OMNeT++.} Simu5G provides end-to-end simulation of 5G NR
spanning application, transport, and network layers, with support for MEC,
C-V2X, and dual connectivity. However, its PHY relies on stochastic 3GPP
abstractions, so it lacks site-specific 3D channel modelling and flexible
MIMO configurations. Most critically for ISAC, the OMNeT++ architecture makes
integrating GPU-accelerated ray tracers such as Sionna RT onerous: the
co-simulation bridge requires heavy glue code, Simu5G exposes only
omnidirectional antennas, and the rich channel outputs computed in Sionna
(per-path gains, delays, AoA/AoD) cannot be propagated cleanly to the upper
layers.

\textbf{5G LENA on ns-3.} 5G LENA implements PHY/MAC system-level behaviour
aligned with 3GPP standards and supports full-stack studies, from packet
capture and IP transport up to RAN-layer scheduling and application
flows~\cite{ref:5glena}. By default it uses simplified stochastic channel
models (e.g., 3GPP TR~38.901~\cite{ref:tr38901}) for path loss, fading, and
beamforming, and therefore lacks deterministic 3D geometry awareness. However,
recent work on Ns3Sionna~\cite{ref:ns3sionna} shows that Sionna RT's
GPU-accelerated ray tracing can be integrated directly into ns-3, yielding
high-fidelity site-specific channels within the existing ns-3 stack; the
published Wi-Fi integration serves as a proof of concept and a solid
technical foundation to extend the same approach to 5G LENA. In this
ns-3/LENA path, integration is cleaner and more expressive: both sides expose
MIMO/beamforming configuration, the per-path channel state computed in Sionna
(complex gains, delays, Doppler, AoA/AoD) can be passed through intact, and
the resulting scheduling, link adaptation, and application-layer metrics
faithfully reflect the underlying geometry-driven radio conditions.

For these reasons, the framework presented in this paper pairs 5G LENA with
Sionna~RT~\cite{ref:sionna_rt} and gives each tool a clear role: Sionna
handles the physical and link layers (ray tracing, MIMO arrays/beamforming,
per-path/beam loss and delay), while 5G LENA consumes these geometry-aware
results to drive the upper layers (RLC/PDCP/IP/applications), so that
buffering, congestion control, and application QoE evolve in response to the
same radio conditions produced by Sionna. This division of labour is detailed
in Section~\ref{sec:methodology}.
